\section{Overview and Scope}
\label{chap:overview}

\subsection{The Problem}

The SAP OSS training pipeline requires high-quality Text-to-SQL training data that reflects
the actual schema structure of SAP HANA Cloud tables. The current implementation in
\texttt{src/training/pipeline/} uses metadata from \textbf{Excel files}
(\texttt{DATA\_DICTIONARY.xlsx}, \texttt{ESG\_DATA\_DICTIONARY.xlsx}) that were originally
exported from HANA or defined manually.

While the \texttt{HanaClient} in \texttt{src/training/pipeline/hana\_client.py} has the
capability to query \texttt{SYS.TABLE\_COLUMNS} for live schema metadata, this HANA-direct
extraction is \textbf{not currently integrated} into the main training data generation
pipeline.

\begin{table}[H]
\centering
\caption{Current vs. Target Training Data Pipeline}
\label{tab:gap-analysis}
\begin{tabularx}{\textwidth}{lXX}
\toprule
\textbf{Dimension} & \textbf{Current State} & \textbf{Target State} \\
\midrule
Schema Source & Static Excel/CSV files & Live HANA Cloud metadata \\
Extraction Method & Manual export & MCP-based automated query \\
Training Data & Template-based expansion & Reasoning-driven synthesis \\
Diversity Control & None & Taxonomic coverage \\
Quality Assurance & None & Double-critic filtering \\
Complexity Control & None & Calibrated Elo scoring \\
\bottomrule
\end{tabularx}
\end{table}

\subsection{Goals}

The specification targets three outcomes:

\begin{enumerate}[nosep]
    \item \textbf{HANA-Direct Extraction.} Every training dataset is generated from live
          schema metadata queried via the AI Core PAL MCP service.
    \item \textbf{Taxonomic Coverage.} The conceptual space of the target domain is mapped
          using hierarchical taxonomies, ensuring global diversity.
    \item \textbf{Reasoning-Driven Synthesis.} Each training example is generated from
          first principles using agentic refinement, ensuring quality, local diversity,
          and controlled complexity.
\end{enumerate}

\subsection{In-Scope Domains}

\begin{table}[H]
\centering
\caption{Application Domains}
\label{tab:domains}
\begin{tabularx}{\textwidth}{llX}
\toprule
\textbf{Domain} & \textbf{HANA Schema} & \textbf{Training Task} \\
\midrule
Arabic Document Processing & \texttt{AP\_INVOICES}, \texttt{VAT\_CHECKLISTS} & Invoice field extraction, VAT compliance \\
Trial Balance Review & \texttt{TB\_VARIANCES}, \texttt{GL\_BALANCES} & Variance analysis, commentary generation \\
ESG Analytics & \texttt{ESG\_METRICS}, \texttt{FINANCED\_EMISSIONS} & Scope classification, net-zero alignment \\
\bottomrule
\end{tabularx}
\end{table}

\subsection{System Context}

\begin{figure}[H]
\centering
\begin{tikzpicture}[
    node distance=1.2cm and 2cm,
    box/.style={rectangle, draw, rounded corners, minimum width=2.5cm, minimum height=1cm, align=center, font=\small},
    hana/.style={box, fill=hanagreen!20, draw=hanagreen},
    mcp/.style={box, fill=sapblue!20, draw=sapblue},
    llm/.style={box, fill=llmpurple!20, draw=llmpurple},
    output/.style={box, fill=sapgold!20, draw=sapgold},
    arr/.style={-{Stealth[length=3mm]}, thick},
]
    % HANA Layer
    \node[hana] (hana) {SAP HANA\\Cloud};
    
    % MCP Layer
    \node[mcp, right=of hana] (mcp) {AI Core PAL\\MCP :8084};
    
    % Extraction
    \node[box, right=of mcp, fill=saplight] (extract) {Schema\\Extractor};
    
    % Registry
    \node[box, below=of extract, fill=saplight] (registry) {Schema\\Registry};
    
    % LLM
    \node[llm, right=2.5cm of extract] (llm) {vLLM\\TurboQuant\\:9180};
    
    % Taxonomy
    \node[box, below=of llm, fill=saplight] (taxonomy) {Taxonomy\\Builder};
    
    % Generator
    \node[box, below=of taxonomy, fill=saplight] (generator) {Data\\Generator};
    
    % Output
    \node[output, right=of generator] (output) {Training Data\\(.jsonl)};
    
    % Arrows
    \draw[arr] (hana) -- node[above, font=\tiny]{SYS.TABLE\_COLUMNS} (mcp);
    \draw[arr] (mcp) -- node[above, font=\tiny]{execute\_sql} (extract);
    \draw[arr] (extract) -- (registry);
    \draw[arr] (registry) -- (taxonomy);
    \draw[arr, dashed] (llm) -- node[right, font=\tiny]{Best-of-N} (taxonomy);
    \draw[arr] (taxonomy) -- (generator);
    \draw[arr, dashed] (llm) -- node[right, font=\tiny]{Critic} (generator);
    \draw[arr] (generator) -- (output);
    
    % Labels
    \node[below=0.3cm of hana, font=\tiny\color{hanagreen}] {Database};
    \node[below=0.3cm of mcp, font=\tiny\color{sapblue}] {Protocol};
    \node[below=0.3cm of llm, font=\tiny\color{llmpurple}] {LLM Backend};
\end{tikzpicture}
\caption{System context. HANA Cloud (green) provides live schema metadata via the MCP
protocol (blue). The vLLM backend (purple) powers taxonomy expansion and data generation.
Output is JSONL training data (gold).}
\label{fig:system-context}
\end{figure}

\subsection{Research Paper Traceability}

This specification implements the Simula framework from Davidson et al. (2026),
``Reasoning-Driven Synthetic Data Generation and Evaluation,'' published in
Transactions on Machine Learning Research (TMLR). The following table maps
paper sections to specification sections:

\begin{table}[H]
\centering
\caption{Traceability: Research Paper to Implementation}
\label{tab:traceability}
\begin{tabularx}{\textwidth}{p{3.5cm}p{4cm}X}
\toprule
\textbf{Paper Reference} & \textbf{Paper Finding} & \textbf{Spec Implementation} \\
\midrule
Section 2.1, Figure 1 & Taxonomy-based coverage & §\ref{chap:taxonomy}, Algorithm 1 \\
Section 2.2, Algorithm 2 & Agentic data synthesis & §\ref{chap:generator}, Algorithm 2 \\
Section 2.3, Algorithm 3 & Calibrated complexity scoring & §\ref{chap:evaluation}.1, Elo algorithm \\
Section 3.1, Figure 3 & Double-critic validation & §\ref{chap:generator}.5, §\ref{chap:evaluation}.4 \\
Section 3.3, Figure 4 & Additive diversity & §\ref{chap:evaluation}.3, Diversity analyzer \\
Section 4.3 & 83\% saturation ratio & §\ref{chap:evaluation}.5, Gap analysis \\
Figure 7 & Complexity vs teacher strength & §\ref{chap:evaluation}.6.1, Adaptive ratio \\
Table 2 (Appendix B) & Taxonomy quality metrics & §\ref{chap:evaluation}.6.5, Quality gate \\
Appendix E.3 & Elo algorithm details & §\ref{chap:evaluation}.1.1, \texttt{\_update\_elo()} \\
\bottomrule
\end{tabularx}
\end{table}

\subsection{Out of Scope}

The following are \textbf{explicitly} out of scope for this document:

\begin{itemize}[nosep]
    \item Model fine-tuning and evaluation (covered by separate LoRA training specification).
    \item Downstream inference and routing (covered by \texttt{03-prompt-routing.tex}).
    \item Trial Balance business process (covered by \texttt{04-trial-balance-ai-specification.tex}).
    \item Arabic invoice pipeline (covered by the Arabic AP specification).
\end{itemize}

% =============================================================================
% CHAPTER 2: CANONICAL DATA SCHEMA
% =============================================================================
\newpage